\subsection{Minimal Assumptions and Domain of Validity}
  \label{subsec:minimal_assumptions}

  We assume that gravitational phenomenology at galactic and sub-cosmological
  scales can be described by an effective potential $\Phi_{\mathrm{eff}}$,
  used here as a convenient parametrization of kinematic inference rather than
  as a fundamental dynamical field.
  The framework satisfies the following requirements.

  (i) The effective dynamics must reduce to standard Newtonian gravity in
  high-density environments,
  ensuring consistency with Solar System tests.
  (ii) Deviations are allowed only in diffuse regimes,
  characterized by low baryonic density or weak gravitational gradients.
  (iii) The transition is governed by at most one additional universal scale
  parameter $a_\star$,
  with no environment-specific tuning.
  (iv) Locality and isotropy are preserved at the effective level.

  Mathematically, the framework shares the Lagrangian structure of the
  quasi-linear MOND theories proposed by Bekenstein and Milgrom~\cite{Bekenstein1984}.
  However, whereas MOND introduces this modification as a fundamental alteration
  of the gravitational law,
  the present approach interprets it as a saturated effective response.
  The modification reflects a finite capacity of the effective description
  to resolve increasingly dilute relational structure~\cite{Beau2026b}.

  This conceptual shift is essential.
  It allows the phenomenological success of MOND at galactic scales to be
  reinterpreted as a regime of descriptive saturation,
  rather than as a modification of gravity.
  The same mechanism then applies naturally to volumetric density dilution
  in cosmological voids.

  The present framework is therefore not intended as a replacement for MOND
  phenomenology.
  It provides an interpretation of its low-acceleration successes
  as an emergent effective limit within a broader descriptive structure.

  To state the relation precisely: the interpolation structure of
  Eq.~\eqref{eq:modified_poisson} is adopted from the quasi-linear MOND
  class as a phenomenological input, not derived here.
  In the deep galactic regime the predictions coincide with those of QUMOND,
  so rotation-curve fits do not by themselves discriminate between the two
  readings.
  What is distinct, and carries the falsifiable content of this paper, lies
  elsewhere: the operational-closure relation $a_\star \sim \kappa\,c\,H(z)$
  tying the galactic saturation scale to the expansion rate with no new
  degree of freedom (Appendix~\ref{app:cosmo_resolution}), the resulting
  environmental bias of local $H_0$ inference, and the void signatures of
  Section~\ref{sec:voids_hubble}.
  Within the underlying programme, bounded relaxation flux uniquely selects
  a Born--Infeld-type effective action~\cite{Beau2026c}; deriving the
  specific saturation function $\mu(x)$ from that selection is an open task
  and is explicitly not claimed here.
